% ============================================================
%  数学教学文档 · 共用教学环境（第四块：环境与盒子）
%  用法：在 preamble.tex、macros.tex 之后 % ============================================================
%  数学教学文档 · 共用教学环境（第四块：环境与盒子）
%  用法：在 preamble.tex、macros.tex 之后 % ============================================================
%  数学教学文档 · 共用教学环境（第四块：环境与盒子）
%  用法：在 preamble.tex、macros.tex 之后 % ============================================================
%  数学教学文档 · 共用教学环境（第四块：环境与盒子）
%  用法：在 preamble.tex、macros.tex 之后 \input{../common/envs}
%  提供：derivation / summary / commonerror / method（讲义）
%        conccard（结论卡片）
%        notezone / noteblank（笔记）
%        mistake / wrong / correct / variant / reflect（错题本）
%        fastcard（速查表）
%        reflection（教学反思）
% ============================================================

\usepackage{tcolorbox}
\tcbuselibrary{skins, breakable}

% ---------- 讲义用 ----------
% 推导：左侧竖线 + 白色标题条
\newtcolorbox{derivation}{
  breakable, enhanced, sharp corners,
  colback=main!3, colframe=main!3,
  boxrule=0pt, leftrule=2pt, left=10pt, top=6pt, bottom=4pt,
  title={\bfseries\color{main!80!black}推导},
  coltitle=main!80!black, fonttitle=\small,
  attach boxed title to top left={xshift=0mm, yshift=-1.5mm},
  boxed title style={sharp corners, colback=white,
    colframe=main!60!black, boxrule=0.4pt},
}

% 本节小结
\newtcolorbox{summary}{
  breakable, enhanced, sharp corners,
  colback=third!5, colframe=third!60!black, boxrule=0.5pt,
  title={\bfseries 本节小结}, coltitle=black, fonttitle=\small,
  top=6pt,
}

% 易错点
\newtcolorbox{commonerror}[1]{
  breakable, enhanced, sharp corners,
  colback=second!4, colframe=second!65!black, boxrule=0.5pt,
  title={\bfseries 易错点：#1}, coltitle=black, fonttitle=\small,
  top=6pt,
}

% 方法归纳
\newtcolorbox{method}{
  breakable, enhanced, sharp corners,
  colback=main!3, colframe=main!55!black, boxrule=0.4pt,
  title={\bfseries 方法归纳}, coltitle=black, fonttitle=\small,
  top=6pt,
}

% ---------- 结论卡片 ----------
% \begin{conccard}[层级]{结论}{适用条件}{易错点} 推导正文 \end{conccard}
\NewDocumentEnvironment{conccard}{ O{二级} m m m +b }{%
  \begin{tcolorbox}[
    breakable, enhanced, sharp corners,
    colback=main!4, colframe=main!65!black, boxrule=0.6pt,
    title={\bfseries #2}, coltitle=black, fonttitle=\small,
    attach boxed title to top left={xshift=2mm, yshift=-2mm},
    boxed title style={sharp corners, colback=white,
      colframe=main!65!black, boxrule=0.4pt},
    top=7pt,
  ]%
  {\small #5}\par\vspace{0.6ex}%
  {\footnotesize
    \textbf{\color{main}适用条件：}#3\par
    \vspace{0.2ex}%
    \textbf{\color{second}易错点：}#4\par
    \vspace{0.2ex}%
    \textbf{\color{third}层级：}#1%
  }%
  \end{tcolorbox}%
}{}

% ---------- 笔记用 ----------
\newtcolorbox{notezone}[1]{
  breakable, enhanced, sharp corners,
  colback=yellow!4, colframe=yellow!45!orange, boxrule=0.5pt,
  title={\bfseries 笔记：#1}, coltitle=black, fonttitle=\small,
  top=6pt,
}
% 留白（供学生手写）
\newcommand{\noteblank}[1][6ex]{\vspace{#1}}

% ---------- 错题本用 ----------
% \begin{mistake}[错因分类] 原题与四个小段 \end{mistake}
\NewDocumentEnvironment{mistake}{ O{概念不清} +b }{%
  \begin{tcolorbox}[
    breakable, enhanced, sharp corners,
    colback=red!3, colframe=red!50!black, boxrule=0.5pt,
    title={\bfseries 错题（#1）}, coltitle=black, fonttitle=\small,
    top=6pt,
  ]#2\end{tcolorbox}%
}{}
\newcommand{\wrong}[1]{\par{\bfseries\color{red!70!black}错解：}#1\par}
\newcommand{\correct}[1]{\par{\bfseries\color{main}正解：}#1\par}
\newcommand{\variant}[1]{\par{\bfseries\color{third}变式：}#1\par}
\newcommand{\reflect}[1]{\par{\bfseries\color{second}反思：}#1\par}

% ---------- 速查表用 ----------
\newtcolorbox{fastcard}[1]{
  enhanced, sharp corners, breakable=false,
  colback=main!3, colframe=main!55!black, boxrule=0.4pt,
  title={\bfseries #1}, coltitle=black, fonttitle=\small,
  top=5pt,
}

% ---------- 答案与解析用 ----------
\newtcolorbox{ansblock}[1]{
  breakable, enhanced, sharp corners,
  colback=green!3, colframe=green!45!black, boxrule=0.4pt,
  title={\bfseries 答案：#1}, coltitle=black, fonttitle=\small,
  top=5pt,
}

% ---------- 教学反思用 ----------
\newtcolorbox{teachreflect}[1]{
  breakable, enhanced, sharp corners,
  colback=third!4, colframe=third!60!black, boxrule=0.4pt,
  title={\bfseries 教学反思：#1}, coltitle=black, fonttitle=\small,
  top=5pt,
}

%  提供：derivation / summary / commonerror / method（讲义）
%        conccard（结论卡片）
%        notezone / noteblank（笔记）
%        mistake / wrong / correct / variant / reflect（错题本）
%        fastcard（速查表）
%        reflection（教学反思）
% ============================================================

\usepackage{tcolorbox}
\tcbuselibrary{skins, breakable}

% ---------- 讲义用 ----------
% 推导：左侧竖线 + 白色标题条
\newtcolorbox{derivation}{
  breakable, enhanced, sharp corners,
  colback=main!3, colframe=main!3,
  boxrule=0pt, leftrule=2pt, left=10pt, top=6pt, bottom=4pt,
  title={\bfseries\color{main!80!black}推导},
  coltitle=main!80!black, fonttitle=\small,
  attach boxed title to top left={xshift=0mm, yshift=-1.5mm},
  boxed title style={sharp corners, colback=white,
    colframe=main!60!black, boxrule=0.4pt},
}

% 本节小结
\newtcolorbox{summary}{
  breakable, enhanced, sharp corners,
  colback=third!5, colframe=third!60!black, boxrule=0.5pt,
  title={\bfseries 本节小结}, coltitle=black, fonttitle=\small,
  top=6pt,
}

% 易错点
\newtcolorbox{commonerror}[1]{
  breakable, enhanced, sharp corners,
  colback=second!4, colframe=second!65!black, boxrule=0.5pt,
  title={\bfseries 易错点：#1}, coltitle=black, fonttitle=\small,
  top=6pt,
}

% 方法归纳
\newtcolorbox{method}{
  breakable, enhanced, sharp corners,
  colback=main!3, colframe=main!55!black, boxrule=0.4pt,
  title={\bfseries 方法归纳}, coltitle=black, fonttitle=\small,
  top=6pt,
}

% ---------- 结论卡片 ----------
% \begin{conccard}[层级]{结论}{适用条件}{易错点} 推导正文 \end{conccard}
\NewDocumentEnvironment{conccard}{ O{二级} m m m +b }{%
  \begin{tcolorbox}[
    breakable, enhanced, sharp corners,
    colback=main!4, colframe=main!65!black, boxrule=0.6pt,
    title={\bfseries #2}, coltitle=black, fonttitle=\small,
    attach boxed title to top left={xshift=2mm, yshift=-2mm},
    boxed title style={sharp corners, colback=white,
      colframe=main!65!black, boxrule=0.4pt},
    top=7pt,
  ]%
  {\small #5}\par\vspace{0.6ex}%
  {\footnotesize
    \textbf{\color{main}适用条件：}#3\par
    \vspace{0.2ex}%
    \textbf{\color{second}易错点：}#4\par
    \vspace{0.2ex}%
    \textbf{\color{third}层级：}#1%
  }%
  \end{tcolorbox}%
}{}

% ---------- 笔记用 ----------
\newtcolorbox{notezone}[1]{
  breakable, enhanced, sharp corners,
  colback=yellow!4, colframe=yellow!45!orange, boxrule=0.5pt,
  title={\bfseries 笔记：#1}, coltitle=black, fonttitle=\small,
  top=6pt,
}
% 留白（供学生手写）
\newcommand{\noteblank}[1][6ex]{\vspace{#1}}

% ---------- 错题本用 ----------
% \begin{mistake}[错因分类] 原题与四个小段 \end{mistake}
\NewDocumentEnvironment{mistake}{ O{概念不清} +b }{%
  \begin{tcolorbox}[
    breakable, enhanced, sharp corners,
    colback=red!3, colframe=red!50!black, boxrule=0.5pt,
    title={\bfseries 错题（#1）}, coltitle=black, fonttitle=\small,
    top=6pt,
  ]#2\end{tcolorbox}%
}{}
\newcommand{\wrong}[1]{\par{\bfseries\color{red!70!black}错解：}#1\par}
\newcommand{\correct}[1]{\par{\bfseries\color{main}正解：}#1\par}
\newcommand{\variant}[1]{\par{\bfseries\color{third}变式：}#1\par}
\newcommand{\reflect}[1]{\par{\bfseries\color{second}反思：}#1\par}

% ---------- 速查表用 ----------
\newtcolorbox{fastcard}[1]{
  enhanced, sharp corners, breakable=false,
  colback=main!3, colframe=main!55!black, boxrule=0.4pt,
  title={\bfseries #1}, coltitle=black, fonttitle=\small,
  top=5pt,
}

% ---------- 答案与解析用 ----------
\newtcolorbox{ansblock}[1]{
  breakable, enhanced, sharp corners,
  colback=green!3, colframe=green!45!black, boxrule=0.4pt,
  title={\bfseries 答案：#1}, coltitle=black, fonttitle=\small,
  top=5pt,
}

% ---------- 教学反思用 ----------
\newtcolorbox{teachreflect}[1]{
  breakable, enhanced, sharp corners,
  colback=third!4, colframe=third!60!black, boxrule=0.4pt,
  title={\bfseries 教学反思：#1}, coltitle=black, fonttitle=\small,
  top=5pt,
}

%  提供：derivation / summary / commonerror / method（讲义）
%        conccard（结论卡片）
%        notezone / noteblank（笔记）
%        mistake / wrong / correct / variant / reflect（错题本）
%        fastcard（速查表）
%        reflection（教学反思）
% ============================================================

\usepackage{tcolorbox}
\tcbuselibrary{skins, breakable}

% ---------- 讲义用 ----------
% 推导：左侧竖线 + 白色标题条
\newtcolorbox{derivation}{
  breakable, enhanced, sharp corners,
  colback=main!3, colframe=main!3,
  boxrule=0pt, leftrule=2pt, left=10pt, top=6pt, bottom=4pt,
  title={\bfseries\color{main!80!black}推导},
  coltitle=main!80!black, fonttitle=\small,
  attach boxed title to top left={xshift=0mm, yshift=-1.5mm},
  boxed title style={sharp corners, colback=white,
    colframe=main!60!black, boxrule=0.4pt},
}

% 本节小结
\newtcolorbox{summary}{
  breakable, enhanced, sharp corners,
  colback=third!5, colframe=third!60!black, boxrule=0.5pt,
  title={\bfseries 本节小结}, coltitle=black, fonttitle=\small,
  top=6pt,
}

% 易错点
\newtcolorbox{commonerror}[1]{
  breakable, enhanced, sharp corners,
  colback=second!4, colframe=second!65!black, boxrule=0.5pt,
  title={\bfseries 易错点：#1}, coltitle=black, fonttitle=\small,
  top=6pt,
}

% 方法归纳
\newtcolorbox{method}{
  breakable, enhanced, sharp corners,
  colback=main!3, colframe=main!55!black, boxrule=0.4pt,
  title={\bfseries 方法归纳}, coltitle=black, fonttitle=\small,
  top=6pt,
}

% ---------- 结论卡片 ----------
% \begin{conccard}[层级]{结论}{适用条件}{易错点} 推导正文 \end{conccard}
\NewDocumentEnvironment{conccard}{ O{二级} m m m +b }{%
  \begin{tcolorbox}[
    breakable, enhanced, sharp corners,
    colback=main!4, colframe=main!65!black, boxrule=0.6pt,
    title={\bfseries #2}, coltitle=black, fonttitle=\small,
    attach boxed title to top left={xshift=2mm, yshift=-2mm},
    boxed title style={sharp corners, colback=white,
      colframe=main!65!black, boxrule=0.4pt},
    top=7pt,
  ]%
  {\small #5}\par\vspace{0.6ex}%
  {\footnotesize
    \textbf{\color{main}适用条件：}#3\par
    \vspace{0.2ex}%
    \textbf{\color{second}易错点：}#4\par
    \vspace{0.2ex}%
    \textbf{\color{third}层级：}#1%
  }%
  \end{tcolorbox}%
}{}

% ---------- 笔记用 ----------
\newtcolorbox{notezone}[1]{
  breakable, enhanced, sharp corners,
  colback=yellow!4, colframe=yellow!45!orange, boxrule=0.5pt,
  title={\bfseries 笔记：#1}, coltitle=black, fonttitle=\small,
  top=6pt,
}
% 留白（供学生手写）
\newcommand{\noteblank}[1][6ex]{\vspace{#1}}

% ---------- 错题本用 ----------
% \begin{mistake}[错因分类] 原题与四个小段 \end{mistake}
\NewDocumentEnvironment{mistake}{ O{概念不清} +b }{%
  \begin{tcolorbox}[
    breakable, enhanced, sharp corners,
    colback=red!3, colframe=red!50!black, boxrule=0.5pt,
    title={\bfseries 错题（#1）}, coltitle=black, fonttitle=\small,
    top=6pt,
  ]#2\end{tcolorbox}%
}{}
\newcommand{\wrong}[1]{\par{\bfseries\color{red!70!black}错解：}#1\par}
\newcommand{\correct}[1]{\par{\bfseries\color{main}正解：}#1\par}
\newcommand{\variant}[1]{\par{\bfseries\color{third}变式：}#1\par}
\newcommand{\reflect}[1]{\par{\bfseries\color{second}反思：}#1\par}

% ---------- 速查表用 ----------
\newtcolorbox{fastcard}[1]{
  enhanced, sharp corners, breakable=false,
  colback=main!3, colframe=main!55!black, boxrule=0.4pt,
  title={\bfseries #1}, coltitle=black, fonttitle=\small,
  top=5pt,
}

% ---------- 答案与解析用 ----------
\newtcolorbox{ansblock}[1]{
  breakable, enhanced, sharp corners,
  colback=green!3, colframe=green!45!black, boxrule=0.4pt,
  title={\bfseries 答案：#1}, coltitle=black, fonttitle=\small,
  top=5pt,
}

% ---------- 教学反思用 ----------
\newtcolorbox{teachreflect}[1]{
  breakable, enhanced, sharp corners,
  colback=third!4, colframe=third!60!black, boxrule=0.4pt,
  title={\bfseries 教学反思：#1}, coltitle=black, fonttitle=\small,
  top=5pt,
}

%  提供：derivation / summary / commonerror / method（讲义）
%        conccard（结论卡片）
%        notezone / noteblank（笔记）
%        mistake / wrong / correct / variant / reflect（错题本）
%        fastcard（速查表）
%        reflection（教学反思）
% ============================================================

\usepackage{tcolorbox}
\tcbuselibrary{skins, breakable}

% ---------- 讲义用 ----------
% 推导：左侧竖线 + 白色标题条
\newtcolorbox{derivation}{
  breakable, enhanced, sharp corners,
  colback=main!3, colframe=main!3,
  boxrule=0pt, leftrule=2pt, left=10pt, top=6pt, bottom=4pt,
  title={\bfseries\color{main!80!black}推导},
  coltitle=main!80!black, fonttitle=\small,
  attach boxed title to top left={xshift=0mm, yshift=-1.5mm},
  boxed title style={sharp corners, colback=white,
    colframe=main!60!black, boxrule=0.4pt},
}

% 本节小结
\newtcolorbox{summary}{
  breakable, enhanced, sharp corners,
  colback=third!5, colframe=third!60!black, boxrule=0.5pt,
  title={\bfseries 本节小结}, coltitle=black, fonttitle=\small,
  top=6pt,
}

% 易错点
\newtcolorbox{commonerror}[1]{
  breakable, enhanced, sharp corners,
  colback=second!4, colframe=second!65!black, boxrule=0.5pt,
  title={\bfseries 易错点：#1}, coltitle=black, fonttitle=\small,
  top=6pt,
}

% 方法归纳
\newtcolorbox{method}{
  breakable, enhanced, sharp corners,
  colback=main!3, colframe=main!55!black, boxrule=0.4pt,
  title={\bfseries 方法归纳}, coltitle=black, fonttitle=\small,
  top=6pt,
}

% ---------- 结论卡片 ----------
% \begin{conccard}[层级]{结论}{适用条件}{易错点} 推导正文 \end{conccard}
\NewDocumentEnvironment{conccard}{ O{二级} m m m +b }{%
  \begin{tcolorbox}[
    breakable, enhanced, sharp corners,
    colback=main!4, colframe=main!65!black, boxrule=0.6pt,
    title={\bfseries #2}, coltitle=black, fonttitle=\small,
    attach boxed title to top left={xshift=2mm, yshift=-2mm},
    boxed title style={sharp corners, colback=white,
      colframe=main!65!black, boxrule=0.4pt},
    top=7pt,
  ]%
  {\small #5}\par\vspace{0.6ex}%
  {\footnotesize
    \textbf{\color{main}适用条件：}#3\par
    \vspace{0.2ex}%
    \textbf{\color{second}易错点：}#4\par
    \vspace{0.2ex}%
    \textbf{\color{third}层级：}#1%
  }%
  \end{tcolorbox}%
}{}

% ---------- 笔记用 ----------
\newtcolorbox{notezone}[1]{
  breakable, enhanced, sharp corners,
  colback=yellow!4, colframe=yellow!45!orange, boxrule=0.5pt,
  title={\bfseries 笔记：#1}, coltitle=black, fonttitle=\small,
  top=6pt,
}
% 留白（供学生手写）
\newcommand{\noteblank}[1][6ex]{\vspace{#1}}

% ---------- 错题本用 ----------
% \begin{mistake}[错因分类] 原题与四个小段 \end{mistake}
\NewDocumentEnvironment{mistake}{ O{概念不清} +b }{%
  \begin{tcolorbox}[
    breakable, enhanced, sharp corners,
    colback=red!3, colframe=red!50!black, boxrule=0.5pt,
    title={\bfseries 错题（#1）}, coltitle=black, fonttitle=\small,
    top=6pt,
  ]#2\end{tcolorbox}%
}{}
\newcommand{\wrong}[1]{\par{\bfseries\color{red!70!black}错解：}#1\par}
\newcommand{\correct}[1]{\par{\bfseries\color{main}正解：}#1\par}
\newcommand{\variant}[1]{\par{\bfseries\color{third}变式：}#1\par}
\newcommand{\reflect}[1]{\par{\bfseries\color{second}反思：}#1\par}

% ---------- 速查表用 ----------
\newtcolorbox{fastcard}[1]{
  enhanced, sharp corners, breakable=false,
  colback=main!3, colframe=main!55!black, boxrule=0.4pt,
  title={\bfseries #1}, coltitle=black, fonttitle=\small,
  top=5pt,
}

% ---------- 答案与解析用 ----------
\newtcolorbox{ansblock}[1]{
  breakable, enhanced, sharp corners,
  colback=green!3, colframe=green!45!black, boxrule=0.4pt,
  title={\bfseries 答案：#1}, coltitle=black, fonttitle=\small,
  top=5pt,
}

% ---------- 教学反思用 ----------
\newtcolorbox{teachreflect}[1]{
  breakable, enhanced, sharp corners,
  colback=third!4, colframe=third!60!black, boxrule=0.4pt,
  title={\bfseries 教学反思：#1}, coltitle=black, fonttitle=\small,
  top=5pt,
}
